\documentclass[11pt]{article}
\usepackage[margin=1in]{geometry}
\usepackage[T1]{fontenc}
\usepackage[utf8]{inputenc}
\usepackage[french]{babel}
\usepackage{amsmath,amssymb,amsthm}
\usepackage[colorlinks=true,linkcolor=blue,urlcolor=blue]{hyperref}
\newtheorem{problem}{Problème}
\newenvironment{refs}{\par\small\noindent\emph{Références :}\begin{itemize}\setlength\itemsep{0pt}}{\end{itemize}\normalsize}
\title{Hors de la Caverne \\ \Large Analyse complexe}
\author{}
\date{}
\begin{document}
\maketitle
\noindent\emph{Des problèmes, pas de réponses.} Chaque problème ci-dessous est posé
sans solution. Les références indiquent où trouver les connaissances nécessaires.
\medskip


\section*{Lycée}

\begin{problem}[{De Moivre et les racines de l'unité --- Lycée}]
\textbf{CA-01.} \textbf{Problème.}
\begin{enumerate}
\item[(a)] Démontrez la formule de Moivre par récurrence : pour tout
entier $ n \ge 1 $,
$ (\cos\theta + i\sin\theta)^n = \cos n\theta + i\sin n\theta $.
\item[(b)] Servez-vous-en pour trouver toutes les solutions complexes de $ z^n = 1 $, montrez que ce sont
$ n $ points régulièrement espacés sur le cercle unité (les sommets d'un $ n $-gone régulier), et
démontrez que pour $ n \ge 2 $ leur somme vaut 0.
\item[(c)] Extrayez de la trigonométrie réelle à partir de l'arithmétique complexe : établissez
la formule donnant $ \cos 3\theta $ comme polynôme en $ \cos\theta $, et calculez
$ \cos 36^\circ $ exactement à l'aide des racines cinquièmes de l'unité.
\end{enumerate}

\end{problem}

\noindent Abraham de Moivre a utilisé la formule dès 1707 dans ses travaux sur les
équations, et Roger Cotes (1714) est passé à un cheveu de la forme exponentielle qu'Euler a
finalement écrite $ e^{i\theta} = \cos\theta + i\sin\theta $. Les racines de l'unité sont le point
de rencontre de l'algèbre et de la géométrie : la construction du 17-gone régulier par Gauss en
1796, à dix-neuf ans --- la découverte qui lui a fait choisir les mathématiques ---, c'est au fond la
question (c) poussée jusqu'à $ n = 17 $.

\begin{refs}
\item Contexte : Formule de Moivre --- Wikipédia (\url{https://fr.wikipedia.org/wiki/Formule_de_Moivre})
\item Contexte : Racine de l'unité --- Wikipédia (\url{https://fr.wikipedia.org/wiki/Racine_de_l'unit%C3%A9})
\item Lecture : Needham, \textit{Visual Complex Analysis}, ch. 1
\end{refs}

\bigskip

\begin{problem}[{Multiplier, c'est faire tourner --- Lycée}]
\textbf{CA-02.} \textbf{Problème.} Représentez $ z = a + bi $ par le point $ (a, b) $ du
plan, de module $ |z| = \sqrt{a^2 + b^2} $ et d'argument l'angle avec le demi-axe réel
positif.
\begin{enumerate}
\item[(a)] Démontrez que $ |zw| = |z|\,|w| $ et que les arguments s'ajoutent par multiplication,
de sorte que multiplier par $ z $ fait tourner le plan de $ \arg z $ et le dilate de $ |z| $.
\item[(b)] Montrez que l'identité $ |zw| = |z|\,|w| $, écrite en coordonnées, est exactement
l'identité classique des deux carrés :
$ (a^2 + b^2)(c^2 + d^2) = (ac - bd)^2 + (ad + bc)^2 $.
\item[(c)] Démontrez l'inégalité triangulaire $ |z + w| \le |z| + |w| $ et déterminez exactement quand
il y a égalité.
\end{enumerate}

\end{problem}

\noindent Les nombres complexes sont entrés en mathématiques par la bande : Cardan et
Bombelli (1545--1572) en avaient besoin pour résoudre les équations du troisième degré à trois
racines réelles. Ils ne sont devenus respectables qu'avec le foyer géométrique que leur ont donné
Wessel (1797), Argand (1806) et Gauss. La question (b) montre que la géométrie se cachait à la vue
de tous depuis des siècles --- l'identité des deux carrés était connue de Diophante et de Brahmagupta
mille ans avant que quiconque ne dessine un nombre complexe.

\begin{refs}
\item Contexte : Plan complexe --- Wikipédia (\url{https://fr.wikipedia.org/wiki/Plan_complexe})
\item Contexte : Identité de Brahmagupta--Fibonacci --- Wikipédia (\url{https://fr.wikipedia.org/wiki/Identit%C3%A9_de_Brahmagupta})
\item Lecture : Needham, \textit{Visual Complex Analysis}, ch. 1
\end{refs}

\bigskip

\begin{problem}[{La règle du parallélogramme --- Lycée, short}]
\textbf{CA-20.} \textbf{Problème.} Démontrez que pour tous nombres complexes $ z $ et $ w $,
\[ |z + w|^2 + |z - w|^2 = 2|z|^2 + 2|w|^2, \]
et expliquez ce que cette identité dit des diagonales et des côtés d'un parallélogramme.
\end{problem}

\noindent La démonstration en une ligne --- développer $ |a|^2 = a\bar{a} $ et regarder les
termes croisés s'annuler --- est une première démonstration que la conjugaison fait le travail que
faisaient autrefois les coordonnées. L'identité est aussi l'empreinte algébrique exacte d'une norme
issue d'un produit scalaire : une norme sur un espace vectoriel réel ou complexe vérifie la règle du
parallélogramme si et seulement si elle est induite par un produit scalaire, ce qui est le théorème
de Jordan--von Neumann de 1935.

\begin{refs}
\item Contexte : Règle du parallélogramme --- Wikipédia (\url{https://fr.wikipedia.org/wiki/R%C3%A8gle_du_parall%C3%A9logramme})
\item Contexte : Nombre complexe --- Wikipédia (\url{https://fr.wikipedia.org/wiki/Nombre_complexe})
\end{refs}

\bigskip

\begin{problem}[{Racines carrées sans trigonométrie --- Lycée, short}]
\textbf{CA-21.} \textbf{Problème.} Trouvez les deux racines carrées complexes de $ 3 + 4i $ en
résolvant $ (x + iy)^2 = 3 + 4i $ avec $ x, y $ réels, sans utiliser la trigonométrie. Démontrez
ensuite que tout nombre complexe non nul possède exactement deux racines carrées, et donnez-en une
formule en fonction de $ a $, $ b $ et $ \sqrt{a^2 + b^2} $ lorsque le nombre est
$ a + bi $.
\end{problem}

\noindent Comparer parties réelles et parties imaginaires transforme le problème en deux
équations réelles, et le module fournit la troisième relation qui les rend solubles --- petite
illustration de la façon dont l'algèbre complexe tient sa propre comptabilité. Que tout nombre
possède exactement $ n $ racines $ n $-ièmes est le premier indice du théorème fondamental de
l'algèbre (CA-07), et la raison pour laquelle les mathématiciens ont cessé de s'excuser pour
$ \sqrt{-1} $.

\begin{refs}
\item Contexte : Racine carrée --- Wikipédia (\url{https://fr.wikipedia.org/wiki/Racine_carr%C3%A9e})
\item Contexte : Nombre complexe --- Wikipédia (\url{https://fr.wikipedia.org/wiki/Nombre_complexe})
\end{refs}

\bigskip

\begin{problem}[{Les nombres complexes comme moteur de géométrie --- Lycée}]
\textbf{CA-22.} \textbf{Problème.} Identifiez le plan à $ \mathbb{C} $, et notez $ \omega = e^{2\pi i/3} $
une racine cubique primitive de l'unité.
\begin{enumerate}
\item[(a)] Démontrez que la rotation du plan d'angle $ \theta $ autour du point $ p $ est
l'application $ z \mapsto p + e^{i\theta}(z - p) $, et que la composée de deux rotations est une
rotation ou une translation --- déterminez exactement quand chaque cas se produit.
\item[(b)] Démontrez que le triangle de sommets distincts $ a, b, c $ est équilatéral si et seulement si
\[ a^2 + b^2 + c^2 = ab + bc + ca, \]
et montrez que cela équivaut à $ a + \omega b + \omega^2 c = 0 $ ou
$ a + \omega^2 b + \omega c = 0 $, selon l'orientation du triangle.
\item[(c)] Démontrez le théorème de Napoléon : construisez vers l'extérieur des triangles équilatéraux sur
les trois côtés d'un triangle quelconque ; alors les centres de ces trois triangles équilatéraux
sont eux-mêmes les sommets d'un triangle équilatéral.
\item[(d)] Démontrez que le triangle de Napoléon de (c) a le même centre de gravité que le triangle
initial.
\end{enumerate}

\end{problem}

\noindent Dès lors qu'une rotation n'est plus qu'une multiplication par $ e^{i\theta} $,
tout un genre de géométrie plane devient une algèbre qu'un élève patient peut mener sans dessin. Le
théorème de Napoléon est attribué --- sans aucune bonne raison --- à l'empereur ; il apparaît pour la
première fois par écrit en 1825, quatre ans après sa mort, et reste la meilleure publicité pour la
méthode des nombres complexes. La question (b) mérite d'être retenue : une équation sans contenu
géométrique apparent se révèle dire \og{} équilatéral \fg{}, et c'est ce qui donne à la méthode son air de
machine.

\begin{refs}
\item Contexte : Théorème de Napoléon --- Wikipédia (\url{https://fr.wikipedia.org/wiki/Th%C3%A9or%C3%A8me_de_Napol%C3%A9on})
\item Contexte : Racine de l'unité --- Wikipédia (\url{https://fr.wikipedia.org/wiki/Racine_de_l'unit%C3%A9})
\item Lecture : Needham, \textit{Visual Complex Analysis}, ch. 1
\item Lecture : Hahn, \textit{Complex Numbers and Geometry} (MAA, 1994)
\end{refs}

\bigskip

\begin{problem}[{Sommer des cosinus avec une série géométrique --- Lycée, short}]
\textbf{CA-32.} \textbf{Problème.} Pour $ \theta $ non multiple de $ 2\pi $, sommez la série
géométrique $ \sum_{k=0}^{n-1} e^{ik\theta} $ et prenez parties réelle et imaginaire pour obtenir
des formules closes pour $ \sum_{k=0}^{n-1} \cos k\theta $ et $ \sum_{k=0}^{n-1} \sin k\theta $.
Déduisez-en l'identité du noyau de Dirichlet
\[ 1 + 2\sum_{k=1}^{n} \cos k\theta \;=\; \frac{\sin\!\left(\left(n + \tfrac{1}{2}\right)\theta\right)}{\sin(\theta/2)} . \]
\end{problem}

\noindent Des sommes trigonométriques d'apparence redoutable ne sont que des séries
géométriques déguisées dès que l'on lit $ \cos k\theta $ comme la partie réelle de
$ (e^{i\theta})^k $ ; la formule d'Euler transforme une page de chasse aux angles en une ligne
d'algèbre. Les identités étaient connues de Lagrange au dix-huitième siècle, et le noyau de la
seconde formule est celui que Dirichlet a placé au centre de sa démonstration de 1829 de la
convergence de la série de Fourier d'une fonction convenable --- il réapparaît, avec le même visage,
dans l'analyse fonctionnelle de ce site.

\begin{refs}
\item Contexte : Formule d'Euler --- Wikipédia (\url{https://fr.wikipedia.org/wiki/Formule_d'Euler})
\item Contexte : Série géométrique --- Wikipédia (\url{https://fr.wikipedia.org/wiki/S%C3%A9rie_g%C3%A9om%C3%A9trique}), Noyau de Dirichlet --- Wikipédia (\url{https://fr.wikipedia.org/wiki/Noyau_de_Dirichlet})
\item Lecture : Needham, \textit{Visual Complex Analysis}, ch. 1
\end{refs}

\bigskip

\begin{problem}[{Le théorème de Ptolémée en une identité --- Lycée}]
\textbf{CA-33.} \textbf{Problème.} Soient $ a, b, c, d $ quatre nombres complexes quelconques, vus comme les
points $ A, B, C, D $ du plan.
\begin{enumerate}
\item[(a)] Vérifiez en développant les deux membres l'identité
\[ (a - b)(c - d) + (a - d)(b - c) \;=\; (a - c)(b - d) . \]
\item[(b)] Déduisez-en l'inégalité de Ptolémée : pour quatre points quelconques du plan,
$ AC \cdot BD \le AB \cdot CD + AD \cdot BC $, où $ XY $ désigne la distance entre
$ X $ et $ Y $.
\item[(c)] Placez maintenant les quatre points sur le cercle unité, $ a = e^{i\alpha} $, $ b = e^{i\beta} $,
$ c = e^{i\gamma} $, $ d = e^{i\delta} $ avec $ 0 \le \alpha < \beta < \gamma < \delta < 2\pi $.
Démontrez la formule de la corde $ |e^{i\theta} - e^{i\varphi}| = 2\left|\sin\frac{\theta - \varphi}{2}\right| $,
et montrez que (b) devient alors une \textit{égalité} --- le théorème de Ptolémée pour un quadrilatère
inscriptible --- en vérifiant l'identité obtenue entre produits de sinus à l'aide des formules de
transformation de produits en sommes.
\item[(d)] Appliquez (c) à quatre sommets d'un pentagone régulier pour démontrer que le rapport d'une
diagonale à un côté est le nombre d'or $ \varphi = \frac{1 + \sqrt{5}}{2} $. (CA-01(c) atteint la
même constante par une autre route, via les racines cinquièmes de l'unité.)
\end{enumerate}

\end{problem}

\noindent Claude Ptolémée a démontré le théorème dans l'\textit{Almageste} vers 150 de notre
ère, et ce n'était pas un ornement : l'égalité de (c) est précisément la règle d'addition et de
soustraction des cordes, la récurrence avec laquelle il a calculé la table des cordes qui a porté
l'astronomie grecque et islamique pendant quatorze siècles --- de la trigonométrie avant l'existence
des fonctions trigonométriques. La démonstration par les nombres complexes comprime tout cela dans
l'identité scolaire de la question (a), qui ne demande ni dessin ni discussion de cas, alors que la
démonstration synthétique classique exige une construction auxiliaire qu'il faut deviner. La
question (b) montre que l'inégalité vaut pour \textit{quatre points quelconques}, en toute
configuration ; le cercle est exactement le cas d'égalité.

\begin{refs}
\item Contexte : Théorème de Ptolémée --- Wikipédia (\url{https://fr.wikipedia.org/wiki/Th%C3%A9or%C3%A8me_de_Ptol%C3%A9m%C3%A9e})
\item Contexte : Inégalité de Ptolémée --- Wikipédia (\url{https://fr.wikipedia.org/wiki/In%C3%A9galit%C3%A9_de_Ptol%C3%A9m%C3%A9e}), Table des cordes de Ptolémée --- Wikipédia (en anglais) (\url{https://en.wikipedia.org/wiki/Ptolemy%27s_table_of_chords})
\item Lecture : Needham, \textit{Visual Complex Analysis}, ch. 1
\item Lecture : Hahn, \textit{Complex Numbers and Geometry} (MAA, 1994)
\end{refs}

\bigskip


\section*{Licence}

\begin{problem}[{Les équations de Cauchy--Riemann --- Licence}]
\textbf{CA-03.} \textbf{Problème.} Écrivez $ f(x + iy) = u(x, y) + i\,v(x, y) $.
\begin{enumerate}
\item[(a)] Démontrez que si $ f $ est dérivable au sens complexe en un point, alors les dérivées
partielles y vérifient les équations de Cauchy--Riemann :
$ \frac{\partial u}{\partial x} = \frac{\partial v}{\partial y} $ et
$ \frac{\partial u}{\partial y} = -\frac{\partial v}{\partial x} $.
\item[(b)] Démontrez la réciproque sous l'hypothèse que les dérivées partielles existent et sont continues
au voisinage du point.
\item[(c)] Montrez que $ f(z) = \bar{z} $ (la conjugaison complexe) ne vérifie ni l'une ni l'autre,
bien qu'elle soit aussi lisse qu'une application du plan puisse l'être --- et expliquez
géométriquement ce que la dérivabilité complexe exige de plus que la simple régularité (qu'advient-il
des angles ?).
\item[(d)] Démontrez que si $ f $ est holomorphe et si $ u $ est constante, alors $ f $ est constante
sur un domaine connexe.
\end{enumerate}

\end{problem}

\noindent Les équations apparaissent dans les travaux de d'Alembert de 1752 sur
l'écoulement des fluides ainsi que chez Euler, mais ce sont Cauchy (à partir de 1814) et Riemann
(thèse de 1851) qui en ont fait la définition de la discipline. Elles encodent l'âme géométrique de
l'analyse complexe : une application holomorphe est, à l'échelle infinitésimale, une rotation
composée d'une dilatation --- conservant les angles partout où la dérivée est non nulle ---, ce qui
explique pourquoi une équation en une variable complexe fait le travail de deux EDP réelles.

\begin{refs}
\item Contexte : Équations de Cauchy--Riemann --- Wikipédia (\url{https://fr.wikipedia.org/wiki/%C3%89quations_de_Cauchy-Riemann})
\item Lecture : Ahlfors, \textit{Complex Analysis}, ch. 2
\item Cours : MIT OCW 18.112 Functions of a Complex Variable (\url{https://ocw.mit.edu/courses/18-112-functions-of-a-complex-variable-fall-2008/})
\end{refs}

\bigskip

\begin{problem}[{La dérivabilité complexe n'est pas la dérivabilité réelle --- Licence}]
\textbf{CA-04.} \textbf{Problème.}
\begin{enumerate}
\item[(a)] Démontrez que $ f(z) = |z|^2 $ est dérivable au sens
complexe en exactement un point de $ \mathbb{C} $, et holomorphe nulle part.
\item[(b)] Montrez que $ f(x + iy) = \sqrt{|xy|} $ vérifie les équations de Cauchy--Riemann à
l'origine sans y être dérivable au sens complexe --- les équations ponctuelles seules ne suffisent
donc pas.
\item[(c)] Une fonction réelle d'une variable réelle peut être dérivable exactement une fois, ou
exactement sept fois. Démontrez qu'une fonction dérivable au sens complexe sur un ouvert y possède
des dérivées de tout ordre. (Vous pouvez utiliser la formule intégrale de Cauchy de CA-06 ; l'enjeu
est de mesurer la force du contraste.)
\end{enumerate}

\end{problem}

\noindent Cet ensemble d'exemples marque la frontière entre deux mondes. La dérivabilité
réelle est une condition ponctuelle et fragile ; la dérivabilité complexe sur un ouvert est une
condition rigide et globale. Savoir combien peu il faut ajouter aux équations de Cauchy--Riemann pour
forcer l'holomorphie est une question réellement délicate, à laquelle répond le théorème de
Looman--Menchoff (1923--1936) : la continuité jointe aux équations en tout point suffit --- un résultat
dont la démonstration est bien plus ardue que tout ce qu'on trouve dans un premier cours.

\begin{refs}
\item Contexte : Fonction holomorphe --- Wikipédia (\url{https://fr.wikipedia.org/wiki/Fonction_holomorphe})
\item Contexte : Théorème de Looman--Menchoff --- Wikipédia (en anglais) (\url{https://en.wikipedia.org/wiki/Looman%E2%80%93Menchoff_theorem})
\item Lecture : Stein \& Shakarchi, \textit{Complex Analysis} (Princeton Lectures in Analysis II), ch. 1
\end{refs}

\bigskip

\begin{problem}[{Premières intégrales de contour --- Licence}]
\textbf{CA-05.} \textbf{Problème.} Pour une courbe $ \gamma $ et une fonction continue $ f $, l'intégrale
de contour $ \int_\gamma f(z)\,dz $ est définie par paramétrage.
\begin{enumerate}
\item[(a)] Calculez $ \int z^n\,dz $ sur le cercle unité parcouru une fois dans le sens direct, pour tout
entier $ n $ (positif, négatif et nul), et observez que le résultat est 0 sauf pour l'unique
valeur $ n = -1 $, où il vaut $ 2\pi i $.
\item[(b)] Démontrez l'inégalité ML :
$ \left| \int_\gamma f\,dz \right| \le \max_\gamma |f| \cdot \operatorname{length}(\gamma) $.
\item[(c)] Montrez que $ \int dz/z $ le long d'un chemin allant de 1 à $ -1 $ dépend du fait que le
chemin passe au-dessus ou au-dessous de l'origine, et expliquez ce que cela dit de la possibilité
d'un logarithme continu sur $ \mathbb{C} \setminus \{0\} $.
\end{enumerate}

\end{problem}

\noindent Cette unique valeur survivante $ 2\pi i $ de la question (a) est sans doute le
calcul le plus lourd de conséquences de toute l'analyse : toute la théorie de Cauchy --- et donc les
résidus, les indices, et la moitié de cette page --- en découle. La question (c) est la découverte qui
a tourmenté la correspondance d'Euler et de d'Alembert sur les logarithmes des nombres négatifs dans
les années 1740 : le logarithme complexe est inévitablement multiforme, première ombre de la
topologie portée sur l'analyse.

\begin{refs}
\item Contexte : Intégrale de contour --- Wikipédia (\url{https://fr.wikipedia.org/wiki/M%C3%A9thodes_de_calcul_d'int%C3%A9grales_de_contour})
\item Lecture : Ahlfors, \textit{Complex Analysis}, ch. 4
\item Cours : MIT OCW 18.112 Functions of a Complex Variable (\url{https://ocw.mit.edu/courses/18-112-functions-of-a-complex-variable-fall-2008/})
\end{refs}

\bigskip

\begin{problem}[{Le théorème de Cauchy et la formule intégrale --- Licence}]
\textbf{CA-06.} \textbf{Problème.}
\begin{enumerate}
\item[(a)] Démontrez le théorème intégral de Cauchy pour un triangle
(argument de Goursat, par quadrisections répétées) : si $ f $ est holomorphe sur un ouvert
contenant un triangle plein, l'intégrale de $ f $ le long de son bord est nulle. Étendez aux
domaines convexes.
\item[(b)] Déduisez-en la formule intégrale de Cauchy : pour $ f $ holomorphe dans un disque et
$ z $ intérieur à un cercle $ C $ contenu dans ce disque,
\[ f(z) = \frac{1}{2\pi i} \int_C \frac{f(w)}{w - z}\,dw. \]
(c) Récoltez les conséquences : $ f $ possède des dérivées de tout ordre, obtenues en dérivant
sous l'intégrale ; $ f $ est localement égale à sa série de Taylor ; et les estimations de Cauchy
$ |f^{(n)}(0)| \le n! \, \max_C |f| \, / \, r^n $ valent sur un cercle de rayon $ r $.
\end{enumerate}

\end{problem}

\noindent Cauchy a annoncé le théorème en 1825 et la formule en 1831 ; la démonstration de
Goursat de 1900 a supprimé l'hypothèse de continuité de $ f' $, révélant qu'une unique dérivée
complexe entraîne avec elle l'analyticité. La formule est le mécanisme de la rigidité : les valeurs
de $ f $ à l'intérieur d'un cercle sont une moyenne pondérée de ses valeurs sur le cercle. Tout le
reste de cette page --- Liouville, le principe du maximum, les résidus, Picard --- n'est que cette
formule, itérée avec une ambition croissante.

\begin{refs}
\item Contexte : Théorème intégral de Cauchy --- Wikipédia (\url{https://fr.wikipedia.org/wiki/Th%C3%A9or%C3%A8me_int%C3%A9gral_de_Cauchy})
\item Contexte : Formule intégrale de Cauchy --- Wikipédia (\url{https://fr.wikipedia.org/wiki/Formule_int%C3%A9grale_de_Cauchy})
\item Lecture : Stein \& Shakarchi, \textit{Complex Analysis}, ch. 2
\end{refs}

\bigskip

\begin{problem}[{Le théorème de Liouville et le théorème fondamental de l'algèbre --- Licence}]
\textbf{CA-07.} \textbf{Problème.}
\begin{enumerate}
\item[(a)] Démontrez le théorème de Liouville à partir des estimations de Cauchy :
une fonction holomorphe et bornée sur $ \mathbb{C} $ tout entier est constante.
\item[(b)] Déduisez-en le
théorème fondamental de l'algèbre : tout polynôme non constant à coefficients complexes admet une
racine dans $ \mathbb{C} $ (considérez $ 1/p(z) $).
\item[(c)] Renforcez Liouville : une fonction entière
vérifiant $ |f(z)| \le C(1 + |z|)^k $ est un polynôme de degré au plus $ k $.
\item[(d)] Concluez
de (b), par récurrence, que tout polynôme de degré $ n $ se factorise en exactement $ n $
facteurs linéaires sur $ \mathbb{C} $.
\end{enumerate}

\end{problem}

\noindent D'Alembert a tenté le théorème fondamental en 1746 ; Gauss en a fait le sujet de
sa thèse de doctorat de 1799 et y est revenu trois fois encore, jamais pleinement satisfait. La
démonstration par l'analyse complexe --- le théorème de Liouville a été publié pour la première fois
par Cauchy en 1844, l'année même où Liouville en faisait cours --- est la plus courte jamais trouvée,
et elle mérite qu'on la médite : le théorème porte sur l'algèbre, mais toute démonstration connue
doit y faire entrer en fraude de l'analyse ou de la topologie. Demandez-vous où celle-ci le
fait.

\begin{refs}
\item Contexte : Théorème de Liouville (variable complexe) --- Wikipédia (\url{https://fr.wikipedia.org/wiki/Th%C3%A9or%C3%A8me_de_Liouville_(variable_complexe)})
\item Contexte : Théorème fondamental de l'algèbre --- Wikipédia (\url{https://fr.wikipedia.org/wiki/Th%C3%A9or%C3%A8me_fondamental_de_l'alg%C3%A8bre})
\item Lecture : Stein \& Shakarchi, \textit{Complex Analysis}, ch. 2
\end{refs}

\bigskip

\begin{problem}[{Principe du maximum et lemme de Schwarz --- Licence}]
\textbf{CA-08.} \textbf{Problème.}
\begin{enumerate}
\item[(a)] Démontrez le principe du maximum : si $ f $ est
holomorphe et non constante sur un ouvert connexe, alors $ |f| $ n'a aucun maximum local
intérieur ; par conséquent, sur un domaine borné, si $ f $ est continue jusqu'au bord, le maximum
de $ |f| $ est atteint sur le bord.
\item[(b)] Déduisez-en le lemme de Schwarz : si $ f $ envoie le disque
unité dans lui-même de façon holomorphe avec $ f(0) = 0 $, alors $ |f(z)| \le |z| $ pour tout
$ z $ et $ |f'(0)| \le 1 $, l'égalité en un point quelconque (ou dans la dérivée) forçant
$ f $ à être une rotation.
\item[(c)] Utilisez le lemme de Schwarz pour déterminer toutes les bijections holomorphes du disque unité
sur lui-même.
\end{enumerate}

\end{problem}

\noindent Le principe du maximum exprime la moyenne contenue dans la formule de Cauchy :
une moyenne de valeurs ne peut pas dépasser toutes ces valeurs. Le lemme de Schwarz --- issu des cours
berlinois de Hermann Amandus Schwarz vers 1869, distillé sous sa forme et son nom modernes par
Carathéodory --- a l'air d'une estimation d'écolier et il est en fait le germe de la géométrie
hyperbolique en théorie des fonctions : la question (c) calcule le groupe des isométries du plan
hyperbolique sans jamais le nommer. La réinterprétation de Pick en 1916 a rendu cela exact.

\begin{refs}
\item Contexte : Principe du maximum --- Wikipédia (\url{https://fr.wikipedia.org/wiki/Principe_du_maximum})
\item Contexte : Lemme de Schwarz --- Wikipédia (\url{https://fr.wikipedia.org/wiki/Lemme_de_Schwarz})
\item Lecture : Ahlfors, \textit{Complex Analysis}, ch. 4
\end{refs}

\bigskip

\begin{problem}[{Intégrales réelles par les résidus --- Licence}]
\textbf{CA-09.} \textbf{Problème.}
\begin{enumerate}
\item[(a)] Énoncez et démontrez le théorème des résidus pour une fonction
ayant un nombre fini de pôles à l'intérieur d'un lacet simple. Calculez ensuite trois intégrales
réelles classiques par les méthodes de contour, en justifiant chaque passage à la limite :
\item[(b)] $ \int_{-\infty}^{\infty} \frac{dx}{1 + x^2} = \pi $, à l'aide d'un contour semi-circulaire
(vérifiez le résultat avec l'arctangente) ;
\item[(c)] $ \int_0^\infty \frac{\sin x}{x}\,dx = \frac{\pi}{2} $, à l'aide de $ e^{iz}/z $ et d'un
contour échancré autour de l'origine --- notez que cette intégrale ne converge que
conditionnellement, de sorte que les estimations méritent du soin ;
\item[(d)] $ \int_0^{2\pi} \frac{d\theta}{2 + \cos\theta} $, en substituant
$ z = e^{i\theta} $ pour la convertir en une intégrale de contour sur le cercle unité.
\end{enumerate}

\end{problem}

\noindent Cauchy a créé le calcul des résidus dans les années 1820 et l'a utilisé
exactement de cette manière --- pour évaluer des intégrales réelles qui résistent aux méthodes
réelles. La question (c) est l'intégrale de Dirichlet, calculée par Euler en 1781 et centrale dans
la théorie de la convergence des séries de Fourier de Dirichlet en 1829 ; c'est aussi l'exemple
standard d'une intégrale qui existe comme intégrale de Riemann impropre mais non comme intégrale de
Lebesgue (voir la page d'analyse réelle). Les résidus ont transformé une astuce en industrie : les
tables d'intégrales sont, pour une large part, des arguments de contour congelés.

\begin{refs}
\item Contexte : Théorème des résidus --- Wikipédia (\url{https://fr.wikipedia.org/wiki/Th%C3%A9or%C3%A8me_des_r%C3%A9sidus})
\item Contexte : Intégrale de Dirichlet --- Wikipédia (\url{https://fr.wikipedia.org/wiki/Int%C3%A9grale_de_Dirichlet})
\item Lecture : Ahlfors, \textit{Complex Analysis}, ch. 4
\item Cours : MIT OCW 18.112 Functions of a Complex Variable (\url{https://ocw.mit.edu/courses/18-112-functions-of-a-complex-variable-fall-2008/})
\end{refs}

\bigskip

\begin{problem}[{Compter les zéros avec Rouché --- Licence}]
\textbf{CA-10.} \textbf{Problème.}
\begin{enumerate}
\item[(a)] Démontrez le principe de l'argument : pour $ f $ méromorphe
à l'intérieur et sur un lacet simple $ \gamma $ sans zéro ni pôle sur $ \gamma $,
$ \frac{1}{2\pi i} \int_\gamma \frac{f'}{f}\,dz $ vaut le nombre de zéros moins le
nombre de pôles à l'intérieur, comptés avec multiplicité.
\item[(b)] Déduisez-en le théorème de Rouché : si
$ |g| < |f| $ sur $ \gamma $, alors $ f $ et $ f + g $ ont le même nombre de zéros
à l'intérieur de $ \gamma $.
\item[(c)] Mettez-le au travail : déterminez combien de zéros
$ p(z) = z^4 + 6z + 3 $ possède dans le disque $ |z| < 1 $, et combien dans la couronne
$ 1 < |z| < 2 $.
\item[(d)] Donnez une seconde démonstration du théorème fondamental de l'algèbre, longue
d'une ligne, en comparant un polynôme à son terme dominant sur un grand cercle.
\end{enumerate}

\end{problem}

\noindent Eugène Rouché a publié le théorème en 1862 ; il convertit le comptage des
zéros --- une question d'existence --- en une vérification d'inégalité sur un bord, l'un des plus purs
exemples de la philosophie du principe de l'argument selon laquelle les fonctions analytiques sont
contrôlées par leur comportement au bord. Le théorème reste un outil quotidien : les numériciens
s'en servent pour certifier la localisation des racines, et la théorie de la stabilité en ingénierie
est bâtie sur des arguments exactement de cette forme.

\begin{refs}
\item Contexte : Théorème de Rouché --- Wikipédia (\url{https://fr.wikipedia.org/wiki/Th%C3%A9or%C3%A8me_de_Rouch%C3%A9})
\item Contexte : Principe de l'argument --- Wikipédia (\url{https://fr.wikipedia.org/wiki/Principe_de_l'argument})
\item Lecture : Stein \& Shakarchi, \textit{Complex Analysis}, ch. 3
\end{refs}

\bigskip

\begin{problem}[{Applications conformes des domaines standard --- Licence}]
\textbf{CA-11.} \textbf{Problème.} Une application conforme est une bijection holomorphe entre
domaines.
\begin{enumerate}
\item[(a)] Vérifiez que la transformation de Cayley $ w = \frac{z - i}{z + i} $ envoie conformément le
demi-plan supérieur sur le disque unité, et trouvez sa réciproque.
\item[(b)] Montrez que $ w = e^z $
envoie la bande horizontale $ 0 < \operatorname{Im} z < \pi $ sur le demi-plan supérieur,
et trouvez une application conforme de la bande sur le disque.
\item[(c)] Envoyez le quart de plan
$ \{ \operatorname{Re} z > 0,\ \operatorname{Im} z > 0 \} $ sur le disque, et le
demi-disque supérieur sur le demi-plan supérieur.
\item[(d)] Démontrez que les transformations de Möbius
$ z \mapsto \frac{az + b}{cz + d} $ envoient cercles-et-droites sur cercles-et-droites, et servez-vous
de la transformation de Cayley pour résoudre un problème de Dirichlet : trouvez la fonction
harmonique bornée sur le disque valant 1 au bord sur le demi-cercle supérieur et 0 sur
l'inférieur.
\end{enumerate}

\end{problem}

\noindent La thèse de Riemann de 1851 a refondu l'analyse complexe en géométrie : ce qui
compte dans un domaine, c'est sa forme à équivalence conforme près. Les applications standard de ce
problème constituent le manuel de conversation de l'analyste au travail, et la question (d) montre
pourquoi les ingénieurs les ont apprises --- les applications conformes transportent les fonctions
harmoniques, si bien que l'électrostatique, l'écoulement de la chaleur et la conception de profils
d'aile dans des régions biscornues se ramènent au disque. Schwarz et Christoffel (1867--1869) ont
étendu le dictionnaire à tous les polygones.

\begin{refs}
\item Contexte : Transformation conforme --- Wikipédia (\url{https://fr.wikipedia.org/wiki/Transformation_conforme})
\item Contexte : Transformation de Möbius --- Wikipédia (\url{https://fr.wikipedia.org/wiki/Transformation_de_M%C3%B6bius})
\item Lecture : Needham, \textit{Visual Complex Analysis}
\item Lecture : Ahlfors, \textit{Complex Analysis}, ch. 6
\end{refs}

\bigskip

\begin{problem}[{Les zéros sont isolés : le théorème d'identité --- Licence}]
\textbf{CA-12.} \textbf{Problème.}
\begin{enumerate}
\item[(a)] Démontrez que si $ f $ est holomorphe et non
identiquement nulle sur un ouvert connexe, alors chaque zéro de $ f $ est isolé et d'ordre fini :
au voisinage d'un zéro $ a $, $ f(z) = (z - a)^m g(z) $ avec $ g(a) \ne 0 $.
\item[(b)] Déduisez-en le
théorème d'identité : deux fonctions holomorphes sur un ouvert connexe qui coïncident sur un
ensemble ayant un point d'accumulation dans cet ouvert coïncident partout.
\item[(c)] Concluez que les identités familières persistent
automatiquement : $ \sin^2 z + \cos^2 z = 1 $ pour tout $ z $ complexe, parce que cela vaut sur
$ \mathbb{R} $.
\item[(d)] Montrez par un exemple que (b) échoue de façon spectaculaire pour les fonctions réelles
indéfiniment dérivables (trouvez une fonction lisse non nulle s'annulant sur toute une
demi-droite).
\end{enumerate}

\end{problem}

\noindent C'est le point de vue de Weierstrass par les séries entières qui paie : les
fonctions holomorphes se comportent comme des polynômes de degré infini, et la question (d) donne le
sens précis dans lequel la régularité réelle est molle là où l'analyticité complexe est rigide. Le
théorème d'identité est aussi le permis d'exercer du prolongement analytique (CA-14) : une fonction
holomorphe admet au plus un prolongement à un domaine connexe plus grand, si bien que les
prolongements, quand ils existent, sont des découvertes et non des choix.

\begin{refs}
\item Contexte : Théorème d'identité (principe des zéros isolés) --- Wikipédia (en anglais) (\url{https://en.wikipedia.org/wiki/Identity_theorem})
\item Lecture : Stein \& Shakarchi, \textit{Complex Analysis}, ch. 2
\end{refs}

\bigskip

\begin{problem}[{Trouver une conjuguée harmonique --- Licence, short}]
\textbf{CA-23.} \textbf{Problème.} Vérifiez que $ u(x, y) = x^3 - 3xy^2 $ est harmonique sur
$ \mathbb{R}^2 $, et trouvez toutes les conjuguées harmoniques $ v $ --- c'est-à-dire toutes les
$ v $ pour lesquelles $ f = u + iv $ est holomorphe sur $ \mathbb{C} $. Identifiez la fonction
$ f $ obtenue comme fonction élémentaire de $ z $.
\end{problem}

\noindent La recette consiste à intégrer les équations de Cauchy--Riemann (CA-03) l'une
après l'autre, et l'indétermination sur $ v $ est exactement une constante réelle additive sur un
domaine connexe. L'exercice est le visage concret d'un fait structurel : sur un domaine simplement
connexe, toute fonction harmonique réelle est la partie réelle d'une fonction holomorphe, ce qui
explique pourquoi la théorie du potentiel en dimension deux et l'analyse complexe sont une seule et
même discipline.

\begin{refs}
\item Contexte : Conjuguée harmonique --- Wikipédia (en anglais) (\url{https://en.wikipedia.org/wiki/Harmonic_conjugate})
\item Contexte : Fonction harmonique --- Wikipédia (\url{https://fr.wikipedia.org/wiki/Fonction_harmonique})
\end{refs}

\bigskip

\begin{problem}[{Une fonction, deux séries de Laurent --- Licence, short}]
\textbf{CA-24.} \textbf{Problème.} Développez
\[ f(z) = \frac{1}{z(z-1)} \]
en série de Laurent dans chacune des deux couronnes $ 0 < |z| < 1 $ et $ |z| > 1 $.
Déterminez lequel des deux coefficients de $ z^{-1} $ est le résidu de $ f $ en $ 0 $, et
expliquez soigneusement pourquoi l'autre ne l'est pas.
\end{problem}

\noindent Les deux développements sont tous deux corrects et ils diffèrent précisément dans
le coefficient qui compte, ce qui est le plus limpide des avertissements : une série de Laurent
appartient à une couronne, non à une fonction. Le résidu est par définition le coefficient de
$ z^{-1} $ du développement dans un disque épointé autour du point, celui dont la série a le droit
de calculer l'intégrale de contour --- distinction qui ruine silencieusement les calculs de résidus
lorsqu'on l'oublie.

\begin{refs}
\item Contexte : Série de Laurent --- Wikipédia (\url{https://fr.wikipedia.org/wiki/S%C3%A9rie_de_Laurent})
\item Contexte : Théorème des résidus --- Wikipédia (\url{https://fr.wikipedia.org/wiki/Th%C3%A9or%C3%A8me_des_r%C3%A9sidus})
\end{refs}

\bigskip

\begin{problem}[{Birapport et sphère de Riemann --- Licence}]
\textbf{CA-25.} \textbf{Problème.} Soit $ \hat{\mathbb{C}} = \mathbb{C} \cup \{\infty\} $ la sphère de Riemann,
et soit une transformation de Möbius une application $ T(z) = \frac{az + b}{cz + d} $ avec
$ ad - bc \ne 0 $, prolongée à $ \hat{\mathbb{C}} $ de la manière habituelle.
\begin{enumerate}
\item[(a)] Démontrez que les transformations de Möbius forment un groupe agissant de façon strictement
3-transitive sur $ \hat{\mathbb{C}} $ : étant donnés $ z_1, z_2, z_3 $ distincts, il existe
exactement une $ T $ avec $ T(z_1) = 0 $, $ T(z_2) = 1 $, $ T(z_3) = \infty $. Écrivez-la ;
la valeur $ (z, z_1, z_2, z_3) = T(z) $ est le \textit{birapport}.
\item[(b)] Déduisez-en que le birapport est invariant par toute transformation de Möbius, et démontrez que
quatre points distincts de $ \hat{\mathbb{C}} $ sont sur un même cercle ou une même droite si et
seulement si leur birapport est réel. Retrouvez ainsi la conservation des cercles démontrée à la
main en CA-11(d).
\item[(c)] Démontrez que le groupe des transformations de Möbius est isomorphe à
$ \mathrm{PGL}_2(\mathbb{C}) $, et que par projection stéréographique ces applications sont
exactement les bijections holomorphes de la sphère.
\item[(d)] Classifiez les transformations de Möbius autres que l'identité : montrez que chacune a un ou
deux points fixes sur $ \hat{\mathbb{C}} $, et qu'après normalisation $ ad - bc = 1 $ la classe
de conjugaison est déterminée par $ \mathrm{tr}^2 = (a + d)^2 $ --- parabolique lorsque cela vaut
$ 4 $, elliptique lorsque c'est réel dans $ [0, 4) $, hyperbolique lorsque c'est réel et
supérieur à $ 4 $, loxodromique sinon.
\end{enumerate}

\end{problem}

\noindent Le birapport est plus ancien que l'analyse complexe : c'est l'invariant de la
géométrie projective, connu de Pappus dans l'Antiquité et central dans les travaux de Möbius (1827)
et de Chasles. Le lire du point de vue de l'analyse complexe fond deux disciplines --- le groupe de
Möbius est à la fois le groupe des automorphismes de la sphère, la droite projective sur
$ \mathbb{C} $, et (sous les traits de $ \mathrm{PSL}_2(\mathbb{C}) $) le groupe des isométries
de l'espace hyperbolique de dimension 3 et le groupe de Lorentz agissant sur la sphère céleste de la
relativité restreinte. La classification par la trace de la question (d) est ce qui rend possibles
les images de la théorie des groupes discrets de Klein et de Poincaré.

\begin{refs}
\item Contexte : Transformation de Möbius --- Wikipédia (\url{https://fr.wikipedia.org/wiki/Transformation_de_M%C3%B6bius})
\item Contexte : Birapport --- Wikipédia (\url{https://fr.wikipedia.org/wiki/Birapport})
\item Contexte : Sphère de Riemann --- Wikipédia (\url{https://fr.wikipedia.org/wiki/Sph%C3%A8re_de_Riemann})
\item Lecture : Needham, \textit{Visual Complex Analysis}, ch. 3
\item Lecture : Ahlfors, \textit{Complex Analysis}, ch. 3
\end{refs}

\bigskip

\begin{problem}[{Le théorème de l'image ouverte --- Licence, short}]
\textbf{CA-34.} \textbf{Problème.} Démontrez le théorème de l'image ouverte : une fonction
holomorphe non constante sur un ouvert connexe de $ \mathbb{C} $ envoie les ouverts sur des
ouverts. Déduisez-en ensuite en deux lignes le principe du maximum de CA-08(a), et expliquez
pourquoi le théorème est faux pour les applications réelles différentiables du plan et même pour les
fonctions réelles analytiques non constantes d'une variable réelle.
\end{problem}

\noindent La démonstration usuelle passe par le théorème d'application locale : au
voisinage d'un point où $ f - f(z_0) $ a un zéro d'ordre $ m $, la fonction est $ m $-à-un sur
un petit disque, ce que l'on établit avec le théorème de Rouché (CA-10) ou le principe de
l'argument. L'ouverture est le visage géométrique de la rigidité qui traverse cette page --- une image
holomorphe ne peut avoir aucun point de bord provenant de l'intérieur, donc $ |f| $ ne peut avoir
aucun maximum intérieur --- et c'est l'analogue exact, dans un monde complètement différent, du
théorème de l'application ouverte de Banach pour les opérateurs entre espaces de Banach.

\begin{refs}
\item Contexte : Théorème de l'image ouverte (analyse complexe) --- Wikipédia (\url{https://fr.wikipedia.org/wiki/Th%C3%A9or%C3%A8me_de_l'image_ouverte})
\item Contexte : Principe du maximum --- Wikipédia (\url{https://fr.wikipedia.org/wiki/Principe_du_maximum})
\item Lecture : Stein \& Shakarchi, \textit{Complex Analysis}, ch. 8
\end{refs}

\bigskip

\begin{problem}[{Produits de Blaschke --- Licence}]
\textbf{CA-35.} \textbf{Problème.} Pour $ a $ dans le disque unité $ \mathbb{D} $, posez
\[ \varphi_a(z) \;=\; \frac{z - a}{1 - \bar{a}z} . \]
\begin{enumerate}
\item[(a)] Démontrez l'identité
\[ 1 - |\varphi_a(z)|^2 \;=\; \frac{(1 - |a|^2)(1 - |z|^2)}{|1 - \bar{a}z|^2} , \]
et déduisez-en que $ \varphi_a $ envoie $ \mathbb{D} $ sur $ \mathbb{D} $ et le cercle unité
sur le cercle unité, avec $ \varphi_a^{-1} = \varphi_{-a} $.
\item[(b)] Appelez $ B(z) = e^{i\theta} \prod_{j=1}^{n} \varphi_{a_j}(z) $ un \textit{produit de Blaschke fini}
de degré $ n $. Démontrez que $ |B| = 1 $ sur le cercle unité et que $ B $ prend toute valeur
$ w \in \mathbb{D} $ exactement $ n $ fois dans $ \mathbb{D} $, comptées avec
multiplicité.
\item[(c)] Démontrez la réciproque : si $ f $ est holomorphe sur $ \mathbb{D} $, continue sur le disque
fermé, et $ |f(z)| = 1 $ pour tout $ |z| = 1 $, alors $ f $ est un produit de Blaschke fini.
Concluez que les produits de Blaschke de degré un sont exactement les automorphismes trouvés en
CA-08(c).
\item[(d)] Produits infinis : soit $ (a_n) $ une suite de $ \mathbb{D} \setminus \{0\} $ avec
$ \sum_n (1 - |a_n|) < \infty $. Démontrez que
\[ B(z) \;=\; \prod_{n} \frac{|a_n|}{a_n}\, \varphi_{a_n}(z) \]
converge uniformément sur tout compact de $ \mathbb{D} $ vers une fonction holomorphe bornée par 1
dont les zéros sont exactement les $ a_n $. Démontrez ensuite la réciproque à l'aide de la formule
de Jensen (CA-26) : si une fonction holomorphe bornée sur $ \mathbb{D} $, non identiquement nulle,
a pour zéros les $ (a_n) $, alors ces zéros doivent vérifier
$ \sum_n (1 - |a_n|) < \infty $.
\end{enumerate}

\end{problem}

\noindent Wilhelm Blaschke a introduit ces produits et la condition de convergence en 1915,
et la question (d) est la réponse fine à une question que le théorème de factorisation de
Weierstrass (CA-16) laisse ouverte : quelles suites peuvent être l'ensemble des zéros d'une fonction
holomorphe \textit{bornée} sur le disque --- la réponse est une condition sur la vitesse à laquelle les
zéros peuvent approcher le bord. Les produits de Blaschke finis sont les polynômes propres au
disque : ce sont exactement les auto-applications holomorphes propres, et par (c) ils sont
déterminés par leur seul comportement au bord. Les produits infinis sont les premiers exemples de
fonctions \textit{intérieures}, et par le théorème de Beurling (CA-19(a)) les fonctions intérieures
indexent les sous-espaces invariants du décalage --- c'est ainsi que cet innocent facteur de Möbius
finit au centre de la théorie des opérateurs.

\begin{refs}
\item Contexte : Produit de Blaschke --- Wikipédia (en anglais) (\url{https://en.wikipedia.org/wiki/Blaschke_product})
\item Contexte : Espace de Hardy (fonctions intérieures et extérieures) --- Wikipédia (\url{https://fr.wikipedia.org/wiki/Espace_de_Hardy})
\item Lecture : Garnett, \textit{Bounded Analytic Functions}, ch. II
\item Lecture : Conway, \textit{Functions of One Complex Variable I}, ch. VII
\item Cours : MIT OCW 18.112 Functions of a Complex Variable (\url{https://ocw.mit.edu/courses/18-112-functions-of-a-complex-variable-fall-2008/})
\end{refs}

\bigskip


\section*{Master}

\begin{problem}[{Le théorème de l'application conforme de Riemann --- Master}]
\textbf{CA-13.} \textbf{Problème.} (Problème de compréhension, avec une démonstration à l'intérieur.) Le théorème
de l'application conforme de Riemann : tout ouvert simplement connexe de $ \mathbb{C} $ autre que
$ \mathbb{C} $ lui-même est conformément équivalent au disque unité.
\begin{enumerate}
\item[(a)] Démontrez vous-même la moitié \og{} unicité \fg{} :
l'existence étant admise, montrez que l'application est unique dès qu'on impose
$ f(z_0) = 0 $ et $ f'(z_0) > 0 $ (utilisez le lemme de Schwarz de CA-08).
\item[(b)] Démontrez que
$ \mathbb{C} $ lui-même doit être exclu, et que le théorème est faux pour le disque épointé --- où
la simple connexité intervient-elle ?
\item[(c)] Lisez la démonstration moderne et rendez compte de son
moteur : qu'est-ce qu'une famille normale, que dit le théorème de Montel, et comment le fait de
maximiser $ |f'(z_0)| $ sur les applications injectives à valeurs dans le disque fait-il surgir la
surjection à partir de rien ?
\end{enumerate}

\end{problem}

\noindent Riemann a énoncé le théorème dans sa thèse de 1851 avec une démonstration
s'appuyant sur le principe de Dirichlet --- dont Weierstrass a ensuite montré qu'il n'était pas
démontré. Il a fallu un demi-siècle pour combler la lacune (Osgood a donné la première démonstration
complète en 1900 ; l'argument par familles normales aujourd'hui standard est venu de Fejér, Riesz et
Carathéodory dans les années 1910--1920), et la combler a construit une grande partie de l'analyse
moderne, y compris le sauvetage du principe de Dirichlet par Hilbert. Le théorème reste stupéfiant :
l'intérieur d'un flocon fractal déchiqueté et un demi-plan sont le même domaine pour un analyste
complexe.

\begin{refs}
\item Contexte : Théorème de l'application conforme --- Wikipédia (\url{https://fr.wikipedia.org/wiki/Th%C3%A9or%C3%A8me_de_l'application_conforme})
\item Lecture : Ahlfors, \textit{Complex Analysis}, ch. 6
\item Lecture : Stein \& Shakarchi, \textit{Complex Analysis}, ch. 8
\end{refs}

\bigskip

\begin{problem}[{Prolongement analytique et fonction zêta --- Master}]
\textbf{CA-14.} \textbf{Problème.} Pour $ \operatorname{Re} s > 1 $, définissez
$ \zeta(s) = \sum_{n \ge 1} n^{-s} $.
\begin{enumerate}
\item[(a)] Démontrez le produit eulérien :
\[ \zeta(s) = \prod_{p \text{ premier}} \left(1 - p^{-s}\right)^{-1}, \]
et expliquez précisément comment la factorisation unique est utilisée ; déduisez-en la démonstration
par Euler de l'infinité des nombres premiers (que se passe-t-il lorsque $ s \to 1^+ $ ?).
\item[(b)] Démontrez que
$ \zeta(s) - \frac{1}{s - 1} $ se prolonge holomorphiquement à $ \operatorname{Re} s > 0 $,
par exemple via la série alternée $ \eta(s) = \sum (-1)^{n-1} n^{-s} $ ou en comparant la somme à
l'intégrale de $ x^{-s} $.
\item[(c)] Problème de lecture : étudiez le mémoire de Riemann de 1859
(neuf pages) et énoncez soigneusement l'équation fonctionnelle reliant $ \zeta(s) $ à
$ \zeta(1 - s) $, la forme symétrique $ \xi(s) = \xi(1 - s) $, et les rôles exacts de la fonction
Gamma et de la fonction thêta dans la démonstration.
\end{enumerate}

\end{problem}

\noindent Euler a découvert la formule du produit en 1737, faisant des nombres premiers un
objet d'analyse ; l'article de Riemann de 1859 --- son unique travail de théorie des nombres --- a
prolongé $\zeta$ au plan tout entier et lié le terme d'erreur dans la répartition des nombres premiers à la
localisation de ses zéros. Ce problème est la porte d'entrée de l'hypothèse de Riemann (CA-17) : les
questions (a) et (b) sont d'honnêtes exercices, et la question (c) compte parmi les lectures les
plus rentables des mathématiques, neuf pages qui ont créé la théorie analytique des nombres.

\begin{refs}
\item Contexte : Fonction zêta de Riemann --- Wikipédia (\url{https://fr.wikipedia.org/wiki/Fonction_z%C3%AAta_de_Riemann})
\item Contexte : Produit eulérien --- Wikipédia (\url{https://fr.wikipedia.org/wiki/Produit_eul%C3%A9rien})
\item Lecture : Stein \& Shakarchi, \textit{Complex Analysis}, ch. 6
\item Lecture : Titchmarsh, \textit{The Theory of the Riemann Zeta-function}
\end{refs}

\bigskip

\begin{problem}[{Les théorèmes de Picard --- Master}]
\textbf{CA-15.} \textbf{Problème.}
\begin{enumerate}
\item[(a)] Démontrez le théorème de Casorati--Weierstrass : au voisinage
d'une singularité essentielle, une fonction holomorphe s'approche arbitrairement près de toute
valeur complexe.
\item[(b)] Vérifiez Picard sur l'exemple standard : montrez que $ e^z $ omet exactement la valeur 0, et
que $ e^{1/z} $ prend toute valeur non nulle une infinité de fois dans tout voisinage épointé de
0.
\item[(c)] Problème de compréhension : le petit théorème de Picard dit qu'une fonction entière non
constante omet au plus une valeur complexe ; le grand théorème dit qu'au voisinage d'une singularité
essentielle toute valeur, à au plus une exception près, est prise une infinité de fois. Lisez une
démonstration moderne du petit théorème et expliquez le mécanisme qui interdit d'omettre deux
valeurs --- où la géométrie hyperbolique (ou la fonction modulaire, ou le théorème de Bloch, selon
votre source) intervient-elle, et pourquoi une valeur omise est-elle inoffensive alors que deux sont
fatales ?
\end{enumerate}

\end{problem}

\noindent Émile Picard a démontré les deux théorèmes en 1879, et ils ont stupéfié ses
contemporains : Casorati--Weierstrass dit que l'image au voisinage d'une singularité essentielle est
dense, Picard dit qu'elle est tout, sauf un point. Un siècle de redémonstrations --- l'argument
élémentaire de Borel, les familles normales de Montel, la constante de Bloch, la renormalisation de
Zalcman --- a chaque fois ouvert un nouveau chapitre de la théorie des fonctions, et la recherche de
la \og{} bonne \fg{} démonstration est encore une tradition vivante. Les deux valeurs omises 0 et 1 de la
configuration clé sont les mêmes 0 et 1 qui sous-tendent la fonction modulaire $\lambda$.

\begin{refs}
\item Contexte : Théorèmes de Picard --- Wikipédia (\url{https://fr.wikipedia.org/wiki/Th%C3%A9or%C3%A8mes_de_Picard})
\item Contexte : Théorème de Casorati--Weierstrass --- Wikipédia (\url{https://fr.wikipedia.org/wiki/Th%C3%A9or%C3%A8me_de_Weierstrass-Casorati})
\item Lecture : Ahlfors, \textit{Complex Analysis}
\end{refs}

\bigskip

\begin{problem}[{La factorisation de Weierstrass --- Master}]
\textbf{CA-16.} \textbf{Problème.}
\begin{enumerate}
\item[(a)] Démontrez qu'un produit infini $ \prod (1 + a_n) $
avec $ \sum |a_n| < \infty $ converge vers une limite non nulle, et développez les facteurs
élémentaires $ E_k(z) = (1 - z)\exp\!\left(z + \frac{z^2}{2} + \cdots + \frac{z^k}{k}\right) $.
\item[(b)] Démontrez le théorème de factorisation de Weierstrass : toute fonction entière se factorise en
$ z^m e^{g(z)} \prod_n E_{k_n}(z/a_n) $ sur ses zéros $ a_n $ ; réciproquement, toute suite
vérifiant $ |a_n| \to \infty $ est l'ensemble des zéros d'une certaine fonction entière.
\item[(c)] Établissez la pièce classique de
démonstration :
\[ \sin \pi z = \pi z \prod_{n=1}^{\infty} \left( 1 - \frac{z^2}{n^2} \right), \]
et expliquez comment l'identification des coefficients dans cette identité donne l'évaluation
d'Euler $ \sum 1/n^2 = \pi^2/6 $ --- le problème de Bâle, dont l'histoire est racontée dans
L'art de la démonstration (\url{proofs.html}).
\end{enumerate}

\end{problem}

\noindent Euler a factorisé le sinus comme s'il s'agissait d'un polynôme en 1734 et en a
tiré la somme de Bâle, à l'admiration et au malaise de ses contemporains ; Weierstrass (1876) a
fourni la théorie manquante, en montrant exactement quelles prescriptions de zéros sont possibles et
à quel coût exponentiel, et Hadamard (1893) a affiné la factorisation pour les fonctions d'ordre
fini --- la forme utilisée dans l'étude de $\zeta$. C'est le problème où les produits infinis cessent d'être
une magie formelle pour devenir des instruments.

\begin{refs}
\item Contexte : Théorème de factorisation de Weierstrass --- Wikipédia (\url{https://fr.wikipedia.org/wiki/Th%C3%A9or%C3%A8me_de_factorisation_de_Weierstrass})
\item Lecture : Stein \& Shakarchi, \textit{Complex Analysis}, ch. 5
\item Lecture : Ahlfors, \textit{Complex Analysis}, ch. 5
\end{refs}

\bigskip

\begin{problem}[{La formule de Jensen --- Master, short}]
\textbf{CA-26.} \textbf{Problème.} Soit $ f $ holomorphe sur un voisinage du disque fermé
$ |z| \le R $, avec $ f(0) \ne 0 $, sans zéro sur le cercle $ |z| = R $, et de zéros
$ a_1, \dots, a_N $ dans $ |z| < R $ énumérés avec multiplicité. Démontrez la formule de
Jensen :
\[ \log|f(0)| \;=\; \frac{1}{2\pi} \int_0^{2\pi} \log\big|f(Re^{i\theta})\big| \, d\theta \;-\;
\sum_{k=1}^{N} \log \frac{R}{|a_k|}. \]
\end{problem}

\noindent Jensen a publié la formule en 1899, et elle est la charnière entre la taille
d'une fonction holomorphe et le nombre de ses zéros : la moyenne au bord de $ \log|f| $ paie,
terme après terme, pour chaque zéro intérieur. Compter les zéros de cette façon est le fondement de
la théorie de la distribution des valeurs de Nevanlinna et des estimations classiques du nombre de
zéros de la fonction zêta dans un rectangle (CA-14), où elle s'applique presque mot pour mot.

\begin{refs}
\item Contexte : Formule de Jensen --- Wikipédia (\url{https://fr.wikipedia.org/wiki/Formule_de_Jensen})
\item Lecture : Ahlfors, \textit{Complex Analysis}, ch. 5
\end{refs}

\bigskip

\begin{problem}[{Les trois droites de Hadamard --- Master, short}]
\textbf{CA-27.} \textbf{Problème.} Soit $ f $ holomorphe et bornée sur la bande
$ a < \operatorname{Re} z < b $ et continue sur son adhérence, et posez
$ M(x) = \sup_{y \in \mathbb{R}} |f(x + iy)| $. Démontrez que $ \log M(x) $ est une fonction
convexe de $ x $ sur $ [a, b] $.
\end{problem}

\noindent La démonstration est le principe du maximum (CA-08) appliqué à la fonction
auxiliaire $ f(z) \lambda^{z} $ avec $ \lambda $ choisi pour égaliser les deux bords --- un
argument qui vaut d'être appris comme technique, puisque la même astuce démontre les théorèmes de
Phragmén--Lindelöf. L'observation de Hadamard est devenue, entre les mains de Marcel Riesz (1926) et
de Thorin (1938), le théorème d'interpolation complexe qui gouverne la continuité des opérateurs sur
les espaces $ L^p $ ; la convexité dans une bande est l'un des endroits où l'analyse complexe fait
tourner discrètement l'analyse harmonique moderne.

\begin{refs}
\item Contexte : Théorème des trois droites de Hadamard --- Wikipédia (\url{https://fr.wikipedia.org/wiki/Th%C3%A9or%C3%A8me_des_trois_droites_de_Hadamard})
\item Contexte : Principe du maximum --- Wikipédia (\url{https://fr.wikipedia.org/wiki/Principe_du_maximum})
\end{refs}

\bigskip

\begin{problem}[{Fonctions harmoniques, noyau de Poisson et réflexion --- Master}]
\textbf{CA-28.} \textbf{Problème.} Une fonction réelle $ u $ sur un ouvert du plan est \textit{harmonique} si elle
est deux fois continûment dérivable et $ u_{xx} + u_{yy} = 0 $. Pour $ 0 \le r < 1 $, soit
\[ P_r(\theta) \;=\; \frac{1 - r^2}{1 - 2r\cos\theta + r^2} \]
le noyau de Poisson du disque unité.
\begin{enumerate}
\item[(a)] Démontrez qu'une fonction harmonique possède la propriété de la moyenne sur tout cercle dont le
disque fermé est contenu dans le domaine, et déduisez-en le principe du maximum : une fonction
harmonique sur un domaine borné, continue jusqu'au bord, atteint ses extrema sur le bord et est
déterminée par ses valeurs au bord.
\item[(b)] Démontrez que pour $ g $ continue sur le cercle unité l'intégrale de Poisson
\[ u(re^{i\theta}) \;=\; \frac{1}{2\pi} \int_0^{2\pi} P_r(\theta - t)\, g(e^{it})\, dt \]
est harmonique dans le disque et tend vers $ g(e^{i\theta_0}) $ lorsque
$ re^{i\theta} \to e^{i\theta_0} $ --- le problème de Dirichlet sur le disque est donc résolu par
une formule explicite. Montrez réciproquement que toute fonction harmonique dans le disque et
continue sur son adhérence est sa propre intégrale de Poisson.
\item[(c)] Démontrez que sur un domaine simplement connexe toute fonction harmonique est la partie réelle
d'une fonction holomorphe, et montrez que $ \log|z| $ sur $ \mathbb{C} \setminus \{0\} $ n'a pas
de conjuguée harmonique uniforme. Identifiez l'obstruction, et reliez-la à l'inexistence d'un
logarithme continu en CA-05(c).
\item[(d)] Démontrez le principe de réflexion de Schwarz : si $ u $ est harmonique sur la moitié
supérieure d'un disque centré sur l'axe réel, continue jusqu'au diamètre et nulle sur celui-ci,
alors poser $ u(x, -y) = -u(x, y) $ prolonge $ u $ harmoniquement au disque tout entier.
Déduisez-en le principe de réflexion pour les fonctions holomorphes qui sont réelles sur un segment
de l'axe réel.
\end{enumerate}

\end{problem}

\noindent Poisson a écrit son intégrale dans les années 1820 ; le problème de Dirichlet
qu'elle résout --- trouver une fonction harmonique de valeurs au bord prescrites --- a porté toute la
théorie du potentiel du dix-neuvième siècle, de l'électrostatique à la démonstration du théorème de
l'application conforme par Riemann (CA-13). Lorsque Weierstrass a démoli le principe de Dirichlet,
le cas du disque de la question (b) fut l'un des rares morceaux à survivre intacts, et il est devenu
la matière première des réparations ultérieures : la méthode alternée de Schwarz, le balayage de
Poincaré, la méthode sous-harmonique de Perron en 1923. La question (d) est la contribution propre
de Schwarz et le procédé standard pour prolonger une application conforme à travers un arc de
bord.

\begin{refs}
\item Contexte : Fonction harmonique --- Wikipédia (\url{https://fr.wikipedia.org/wiki/Fonction_harmonique})
\item Contexte : Noyau de Poisson --- Wikipédia (\url{https://fr.wikipedia.org/wiki/Noyau_de_Poisson})
\item Contexte : Problème de Dirichlet --- Wikipédia (\url{https://fr.wikipedia.org/wiki/Probl%C3%A8me_de_Dirichlet})
\item Contexte : Principe de réflexion de Schwarz --- Wikipédia (en anglais) (\url{https://en.wikipedia.org/wiki/Schwarz_reflection_principle})
\item Lecture : Ahlfors, \textit{Complex Analysis}, ch. 4 \S{}6
\item Lecture : Conway, \textit{Functions of One Complex Variable I}, ch. X
\end{refs}

\bigskip

\begin{problem}[{Phragmén--Lindelöf dans un secteur --- Master, short}]
\textbf{CA-36.} \textbf{Problème.} Soit $ \alpha \ge 1/2 $, soit
$ S = \{\, z \ne 0 : |\arg z| < \pi/(2\alpha) \,\} $ le secteur ouvert d'ouverture
$ \pi/\alpha $, et soit $ f $ holomorphe sur $ S $ et continue sur son adhérence, avec
$ |f| \le M $ sur les deux demi-droites du bord et $ |f(z)| \le C e^{|z|^{\beta}} $ partout dans
$ S $ pour certaines constantes $ C $ et $ \beta < \alpha $. Démontrez que $ |f| \le M $
partout sur $ S $, et montrez en exhibant $ e^{z} $ sur le demi-plan droit que la conclusion est
fausse lorsque $ \beta = \alpha $.
\end{problem}

\noindent Edvard Phragmén et Ernst Lindelöf ont publié cette famille de théorèmes dans
\textit{Acta Mathematica} en 1908 : le principe du maximum survit sur les domaines non bornés, mais
seulement si l'on interdit à la fonction de croître plus vite que la géométrie du domaine ne le
permet, et la frontière est optimale. La démonstration est l'astuce déjà rencontrée en CA-27 ---
multiplier par $ e^{-\varepsilon z^{\gamma}} $ avec $ \beta < \gamma < \alpha $, appliquer
le principe du maximum ordinaire, puis faire tendre $ \varepsilon \to 0 $. Le principe est le
permis d'exercer standard pour les déplacements de contour de la théorie analytique des nombres et
pour les théorèmes d'interpolation de l'analyse harmonique, où il faut savoir qu'une fonction bornée
sur deux droites est bornée entre elles.

\begin{refs}
\item Contexte : Principe de Phragmén--Lindelöf --- Wikipédia (\url{https://fr.wikipedia.org/wiki/Principe_de_Phragm%C3%A9n%E2%80%93Lindel%C3%B6f})
\item Contexte : Théorème des trois droites de Hadamard --- Wikipédia (\url{https://fr.wikipedia.org/wiki/Th%C3%A9or%C3%A8me_des_trois_droites_de_Hadamard})
\item Lecture : Stein \& Shakarchi, \textit{Complex Analysis}, ch. 4
\end{refs}

\bigskip

\begin{problem}[{Familles normales : Montel et Hurwitz --- Master}]
\textbf{CA-37.} \textbf{Problème.} Dites qu'une famille $ \mathcal{F} $ de fonctions holomorphes sur un ouvert
$ \Omega \subset \mathbb{C} $ est \textit{normale} si toute suite de $ \mathcal{F} $ admet une
sous-suite convergeant uniformément sur tout compact de $ \Omega $.
\begin{enumerate}
\item[(a)] Démontrez le théorème de convergence de Weierstrass : une limite uniforme locale de fonctions
holomorphes est holomorphe, et les dérivées convergent uniformément localement vers la dérivée de la
limite. (Morera et les estimations de Cauchy font tout le travail.)
\item[(b)] Démontrez le théorème de Montel : une famille uniformément bornée sur chaque compact de
$ \Omega $ est normale. (Utilisez les estimations de Cauchy pour obtenir l'équicontinuité, puis
Arzelà--Ascoli et un argument diagonal sur une exhaustion de $ \Omega $ par des compacts.)
\item[(c)] Démontrez le théorème de Hurwitz : si $ f_n \to f $ uniformément localement sur un
$ \Omega $ connexe et si aucune $ f_n $ n'a de zéro, alors ou bien $ f $ est identiquement
nulle, ou bien $ f $ n'a pas de zéro. (Appliquez le principe de l'argument sur un petit cercle et
faites tendre $ n \to \infty $.)
\item[(d)] Déduisez-en qu'une limite uniforme locale de fonctions holomorphes injectives sur un ouvert
connexe est injective ou constante, et servez-vous-en pour combler la lacune de CA-13(c) : la
fonction extrémale obtenue en maximisant $ |f'(z_0)| $ sur les applications injectives de
$ \Omega $ dans le disque est elle-même injective.
\end{enumerate}

\end{problem}

\noindent Paul Montel a introduit les familles normales dans sa thèse de 1907 et a passé
une carrière à montrer que la compacité dans les espaces de fonctions est le moteur caché de la
théorie classique des fonctions : le théorème de l'application conforme (CA-13), les théorèmes de
Picard (CA-15) et les démonstrations modernes du théorème de Bloch (CA-29) deviennent tous des
arguments d'extraction dès que l'on sait quelles familles sont normales. Le théorème de Hurwitz
(1889) est le fait compagnon selon lequel les zéros ne peuvent apparaître ni disparaître à la limite
sans prévenir ; c'est lui qui fait survivre l'univalence --- la propriété qu'a une fonction schlicht
(CA-18) --- au passage à la limite, et donc lui qui rend possible tout ce style de démonstration par
problème extrémal. Le théorème plus profond de Montel, selon lequel une famille omettant deux
valeurs complexes fixées est automatiquement normale, est l'énoncé dont les théorèmes de Picard
découlent en une page.

\begin{refs}
\item Contexte : Théorème de Montel --- Wikipédia (\url{https://fr.wikipedia.org/wiki/Th%C3%A9or%C3%A8me_de_Montel})
\item Contexte : Famille normale --- Wikipédia (\url{https://fr.wikipedia.org/wiki/Famille_normale}), Théorème de Hurwitz (analyse complexe) --- Wikipédia (\url{https://fr.wikipedia.org/wiki/Th%C3%A9or%C3%A8me_de_Hurwitz_sur_les_suites_de_fonctions_holomorphes})
\item Lecture : Conway, \textit{Functions of One Complex Variable I}, ch. VII
\item Lecture : Schiff, \textit{Normal Families} (Springer Universitext, 1993)
\end{refs}

\bigskip

\begin{problem}[{La fonction Gamma --- Master}]
\textbf{CA-38.} \textbf{Problème.} Pour $ \operatorname{Re} s > 0 $, définissez l'intégrale d'Euler
\[ \Gamma(s) \;=\; \int_0^{\infty} t^{s-1} e^{-t}\, dt . \]
\begin{enumerate}
\item[(a)] Démontrez que l'intégrale converge et définit une fonction holomorphe sur le demi-plan
$ \operatorname{Re} s > 0 $, que $ \Gamma(s + 1) = s\,\Gamma(s) $, et que
$ \Gamma(n + 1) = n! $ pour tout entier $ n \ge 0 $.
\item[(b)] Utilisez l'équation fonctionnelle pour prolonger $ \Gamma $ méromorphiquement à
$ \mathbb{C} $ tout entier, et démontrez que le prolongement a des pôles simples exactement en
$ s = 0, -1, -2, \dots $, de résidu $ (-1)^n/n! $ en $ s = -n $.
\item[(c)] Démontrez la formule des compléments d'Euler
\[ \Gamma(s)\,\Gamma(1 - s) \;=\; \frac{\pi}{\sin \pi s} , \]
par exemple en combinant le produit de Weierstrass
$ 1/\Gamma(s) = s e^{\gamma s} \prod_{n \ge 1} (1 + s/n)e^{-s/n} $ avec le produit pour
$ \sin \pi z $ démontré en CA-16(c). Déduisez-en $ \Gamma(1/2) = \sqrt{\pi} $ et que
$ \Gamma $ n'a pas de zéro.
\item[(d)] Démontrez le théorème de Bohr--Mollerup : une fonction $ f : (0, \infty) \to (0, \infty) $ avec
$ f(1) = 1 $, $ f(x+1) = x f(x) $, et $ \log f $ convexe est nécessairement $ \Gamma $
restreinte à $ (0, \infty) $. Expliquez pourquoi une telle condition supplémentaire est
inévitable --- c'est-à-dire pourquoi l'équation fonctionnelle seule ne détermine pas la fonction.
\end{enumerate}

\end{problem}

\noindent Euler a trouvé l'interpolation de la factorielle en 1729, dans sa correspondance
avec Goldbach et Daniel Bernoulli, et Legendre a plus tard imposé la notation décalée $ \Gamma $
dont tout analyste ronchonne depuis. La fonction est la fonction spéciale la plus tenace des
mathématiques --- elle gouverne l'équation fonctionnelle de $ \zeta $ (CA-14), le volume de la boule
de dimension $ n $, l'intégrale bêta et la moitié des constantes de la théorie des probabilités ---
et elle est réellement transcendante en un sens fort : Hölder a démontré en 1887 que $ \Gamma $ ne
vérifie aucune équation différentielle algébrique. Le théorème de Bohr--Mollerup (1922) est la
réponse à la plainte naturelle selon laquelle l'intégrale d'Euler a l'air arbitraire : la
log-convexité jointe à la récurrence détermine complètement la fonction, ce qui a permis à Artin de
bâtir tout son élégant petit livre sur le sujet à partir de cette seule caractérisation.

\begin{refs}
\item Contexte : Fonction gamma --- Wikipédia (\url{https://fr.wikipedia.org/wiki/Fonction_gamma})
\item Contexte : Théorème de Bohr--Mollerup --- Wikipédia (\url{https://fr.wikipedia.org/wiki/Th%C3%A9or%C3%A8me_de_Bohr-Mollerup})
\item Lecture : Artin, \textit{The Gamma Function} (1931 ; traduction anglaise, Dover)
\item Lecture : Stein \& Shakarchi, \textit{Complex Analysis}, ch. 6
\item Lecture : Whittaker \& Watson, \textit{A Course of Modern Analysis}, ch. XII
\end{refs}

\bigskip

\begin{problem}[{Fonctions elliptiques et double périodicité --- Master}]
\textbf{CA-39.} \textbf{Problème.} Fixez $ \omega_1, \omega_2 \in \mathbb{C} $ avec $ \omega_2/\omega_1 $ non
réel, et soit $ \Lambda = \mathbb{Z}\omega_1 + \mathbb{Z}\omega_2 $ le réseau qu'ils engendrent.
Une fonction méromorphe $ f $ sur $ \mathbb{C} $ est \textit{elliptique} pour $ \Lambda $ si
$ f(z + \omega) = f(z) $ pour tout $ \omega \in \Lambda $.
\begin{enumerate}
\item[(a)] Démontrez les trois théorèmes de Liouville sur les fonctions elliptiques : une fonction
elliptique sans pôle est constante ; la somme des résidus dans un parallélogramme fondamental est
nulle ; et une fonction elliptique prend toute valeur de $ \hat{\mathbb{C}} $ le même nombre de
fois dans un parallélogramme fondamental, comptées avec multiplicité. Déduisez-en qu'une fonction
elliptique non constante est d'ordre au moins 2 --- il n'existe pas de fonction elliptique ayant un
unique pôle simple par période.
\item[(b)] Démontrez que la fonction de Weierstrass
\[ \wp(z) \;=\; \frac{1}{z^2} + \sum_{\omega \in \Lambda \setminus \{0\}}
\left( \frac{1}{(z - \omega)^2} - \frac{1}{\omega^2} \right) \]
converge uniformément sur tout compact de $ \mathbb{C} \setminus \Lambda $, que $ \wp $ est
paire et elliptique avec un pôle double en chaque point du réseau, et que $ \wp' $ est elliptique
d'ordre 3. (La convergence repose sur $ \sum_{\omega \ne 0} |\omega|^{-3} < \infty $ ; la
périodicité est plus facile à démontrer d'abord pour $ \wp' $.)
\item[(c)] Démontrez l'équation différentielle
\[ (\wp')^2 \;=\; 4\wp^3 - g_2 \wp - g_3, \qquad
g_2 = 60 \sum_{\omega \ne 0} \omega^{-4}, \quad g_3 = 140 \sum_{\omega \ne 0} \omega^{-6} , \]
en comparant les développements de Laurent à l'origine et en appliquant (a) à la différence.
\item[(d)] Démontrez que toute fonction elliptique pour $ \Lambda $ est une fraction rationnelle en
$ \wp $ et $ \wp' $ --- d'abord les fonctions paires, comme fractions rationnelles en $ \wp $
seule, puis le cas général en séparant parties paire et impaire. Concluez que
$ z \mapsto (\wp(z), \wp'(z)) $ identifie le tore $ \mathbb{C}/\Lambda $ à la cubique projective
$ y^2 = 4x^3 - g_2 x - g_3 $.
\end{enumerate}

\end{problem}

\noindent Les fonctions elliptiques sont nées à l'envers : les analystes du dix-huitième
siècle étudiaient les \textit{intégrales} elliptiques --- la longueur d'arc d'une ellipse et de la
lemniscate ---, et le geste décisif, accompli par Abel en 1827 et Jacobi en 1829 (Gauss l'avait fait
en privé trente ans plus tôt sans rien publier), fut de les inverser : l'encombrante intégrale
multiforme devenait alors une fonction à deux périodes indépendantes, dotée d'une théorie parfaite.
Liouville a fait cours sur la théorie générale en 1844, et le $ \wp $ de Weierstrass l'a réduite à
une unique fonction canonique ; la question (d) énonce que le corps des fonctions elliptiques est
engendré par $ \wp $ et $ \wp' $ avec l'unique relation de (c). Cette relation est le pont vers
la théorie des nombres : la cubique $ y^2 = 4x^3 - g_2x - g_3 $ est une courbe elliptique, la loi
de groupe qu'elle porte est le théorème d'addition pour $ \wp $, et ce dictionnaire entre tores et
cubiques est ce qui porte la modularité et le dernier théorème de Fermat.

\begin{refs}
\item Contexte : Fonction elliptique --- Wikipédia (\url{https://fr.wikipedia.org/wiki/Fonction_elliptique})
\item Contexte : Fonction elliptique de Weierstrass --- Wikipédia (\url{https://fr.wikipedia.org/wiki/Fonction_elliptique_de_Weierstrass}), Couple fondamental de périodes --- Wikipédia (en anglais) (\url{https://en.wikipedia.org/wiki/Fundamental_pair_of_periods})
\item Lecture : Ahlfors, \textit{Complex Analysis}, ch. 7
\item Lecture : Whittaker \& Watson, \textit{A Course of Modern Analysis}, ch. XX
\item Lecture : Silverman, \textit{The Arithmetic of Elliptic Curves}, ch. VI
\end{refs}

\bigskip


\section*{Recherche}

\begin{problem}[{L'hypothèse de Riemann --- Recherche}]
\textbf{CA-17.} \textbf{Problème.} L'hypothèse de Riemann : tout zéro non trivial de
$ \zeta(s) $ a pour partie réelle $ 1/2 $. Avant de consulter quoi que ce soit, faites
honnêtement le travail préparatoire à partir de CA-14 : montrez à partir du produit eulérien que
$ \zeta(s) \ne 0 $ pour $ \operatorname{Re} s > 1 $, et à partir de l'équation fonctionnelle
que les seuls zéros avec $ \operatorname{Re} s < 0 $ sont les zéros triviaux en
$ -2, -4, -6, \dots $, de sorte que tout le mystère habite la bande critique
$ 0 \le \operatorname{Re} s \le 1 $. Démontrez ensuite l'hypothèse.
\end{problem}

\noindent \textbf{Statut : ouvert} depuis 1859 --- huitième problème de Hilbert (1900),
problème du prix du millénaire (2000), et de l'avis commun la question non résolue la plus profonde
des mathématiques. Hardy a démontré en 1914 qu'une infinité de zéros se trouvent sur la droite
critique, et plus de $ 10^{13} $ zéros ont été vérifiés numériquement ; une démonstration ou une
réfutation réorganiserait la théorie des nombres. Elle figure ici parce qu'elle est, au fond, une
question en forme de théorème sur une seule fonction holomorphe. L'histoire complète ---
reformulations, résultats partiels, conséquences --- se trouve dans
Problèmes ouverts (\url{open-problems.html}).

\begin{refs}
\item Contexte : Hypothèse de Riemann --- Wikipédia (\url{https://fr.wikipedia.org/wiki/Hypoth%C3%A8se_de_Riemann})
\item Contexte : Le mémoire de Riemann de 1859 --- Wikipédia (\url{https://fr.wikipedia.org/wiki/Sur_le_nombre_de_nombres_premiers_inf%C3%A9rieurs_%C3%A0_une_taille_donn%C3%A9e})
\item Lecture : Titchmarsh, \textit{The Theory of the Riemann Zeta-function}
\end{refs}

\bigskip

\begin{problem}[{La conjecture de Bieberbach : un géant tombé --- Recherche}]
\textbf{CA-18.} \textbf{Problème.} Dites que $ f $ est \textit{schlicht} si
$ f(z) = z + a_2 z^2 + a_3 z^3 + \cdots $ est holomorphe et injective sur le disque unité.
Conjecture de Bieberbach (1916) : $ |a_n| \le n $ pour tout $ n $, avec égalité seulement pour
les rotations de la fonction de Koebe
$ k(z) = \frac{z}{(1 - z)^2} = z + 2z^2 + 3z^3 + \cdots $.
\begin{enumerate}
\item[(a)] Démontrez le théorème de l'aire pour les fonctions univalentes sur l'extérieur du disque, et
déduisez-en le résultat propre à Bieberbach $ |a_2| \le 2 $.
\item[(b)] Utilisez $ |a_2| \le 2 $ pour démontrer le théorème du quart
de Koebe : l'image du disque par une fonction schlicht contient le disque de rayon
$ 1/4 $.
\item[(c)] Problème de compréhension : lisez l'histoire et la structure de la démonstration
de de Branges --- l'équation d'évolution de Loewner (1923, utilisée pour $ |a_3| \le 3 $), la
réduction à la conjecture de Milin, et l'inégalité d'Askey--Gasper sur les polynômes de Jacobi qui a
clos l'argument --- et rendez compte des raisons pour lesquelles le problème des coefficients a
résisté soixante-neuf ans.
\end{enumerate}

\end{problem}

\noindent \textbf{Statut : résolu.} Ludwig Bieberbach l'a conjecturée dans une note de bas
de page de 1916 ; les progrès sont venus coefficient par coefficient (Loewner $ |a_3| $, 1923 ;
Garabedian--Schiffer $ |a_4| $, 1955) jusqu'à ce que Louis de Branges démontre la conjecture
entière en 1984, publiée dans \textit{Acta Mathematica} en 1985 après que les mathématiciens du
séminaire de Leningrad eurent vérifié et simplifié l'argument. C'est le modèle du problème
\og{} récemment tombé \fg{} : la longue démonstration d'un outsider, l'incrédulité initiale, la vérification
collective minutieuse, et une technique --- l'équation de Loewner --- revenue en gloire deux décennies
plus tard au sein du SLE de Schramm et de deux médailles Fields.

\begin{refs}
\item Contexte : Conjecture de Bieberbach (théorème de de Branges) --- Wikipédia (\url{https://fr.wikipedia.org/wiki/Conjecture_de_Bieberbach})
\item Lecture : Duren, \textit{Univalent Functions}
\item Article : de Branges, \og{} A proof of the Bieberbach conjecture \fg{}, \textit{Acta Mathematica} 154 (1985)
\end{refs}

\bigskip

\begin{problem}[{La théorie des fonctions à la frontière : sous-espaces invariants et compagnie --- Recherche}]
\textbf{CA-19.} \textbf{Problème.} Le problème du sous-espace invariant demande : tout opérateur
linéaire borné sur un espace de Hilbert complexe séparable de dimension infinie possède-t-il un
sous-espace fermé invariant autre que $ \{0\} $ et l'espace tout entier ? L'essentiel de ce que
l'on sait est de l'analyse complexe à une variable déguisée.
\begin{enumerate}
\item[(a)] Échauffement résolu : démontrez le théorème de Beurling
(1949) --- les sous-espaces invariants du décalage $ f(z) \mapsto z f(z) $ sur l'espace de Hardy
$ H^2 $ du disque sont exactement les ensembles $ \varphi H^2 $ pour les fonctions intérieures
$ \varphi $.
\item[(b)] Lisez en quoi la question correspondante pour le décalage de \textit{Bergman} est plus sauvage :
ses sous-espaces invariants peuvent avoir n'importe quel indice, et une description de type
Beurling (Aleman--Richter--Sundberg, 1996) existe mais le treillis est loin d'être compris ---
affinez une partie quelconque de cette description.
\item[(c)] Attaquez une question ouverte voisine d'une variable complexe :
démontrez ou réfutez la conjecture de Brennan (1978), selon laquelle pour toute application conforme
$ \varphi $ d'un domaine simplement connexe $ \Omega $ sur le disque unité, la dérivée vérifie
$ \iint_\Omega |\varphi'|^p \,dA < \infty $ pour tout $ 4/3 < p < 4 $.
\end{enumerate}

\end{problem}

\noindent \textbf{Statut : ouvert.} Pour les espaces de Banach, des contre-exemples existent
--- Enflo (construit dans les années 1970, publié en 1987) et Read (1984), qui a plus tard construit
un opérateur sur $ \ell^1 $ sans aucun sous-espace invariant ---, mais pour l'espace de Hilbert le
problème reste ouvert en 2026 ; une prépublication d'Enflo de 2023 revendiquant une solution
positive n'a pas, à l'heure où ces lignes sont écrites, été acceptée par les pairs, et l'honnêteté
oblige donc à la dire non résolue. Le contraste Hardy/Bergman de (a)--(b) est l'endroit où la théorie
des opérateurs rencontre la théorie classique des fonctions de cette page, et la conjecture de
Brennan en (c) est un problème ouvert purement d'une variable --- élémentaire à énoncer, relié au
spectre universel des moyennes intégrales des applications conformes, et non résolu.

\begin{refs}
\item Contexte : Problème du sous-espace invariant --- Wikipédia (en anglais) (\url{https://en.wikipedia.org/wiki/Invariant_subspace_problem})
\item Contexte : Espace de Hardy --- Wikipédia (\url{https://fr.wikipedia.org/wiki/Espace_de_Hardy})
\item Lecture : Pommerenke, \textit{Boundary Behaviour of Conformal Maps}
\item Article : Enflo, \og{} On the invariant subspace problem in Hilbert spaces \fg{} (2023, prépublication non expertisée), arXiv:2305.15442 (\url{https://arxiv.org/abs/2305.15442})
\end{refs}

\bigskip

\begin{problem}[{La valeur de la constante de Bloch --- Recherche, short}]
\textbf{CA-29.} \textbf{Problème.} Pour $ f $ holomorphe sur le disque unité avec
$ f'(0) = 1 $, soit $ \beta(f) $ la borne supérieure des rayons $ r $ tels qu'une sous-région
du disque soit envoyée par $ f $ bijectivement sur un disque de rayon $ r $ ; le théorème de
Bloch affirme que $ B = \inf_f \beta(f) $ est strictement positif. Déterminez la valeur exacte de
$ B $.
\end{problem}

\noindent \textbf{Statut : ouvert.} Les bornes connues sont
$ 0{,}4332 < B \le 0{,}4719 $, l'inférieure due à Chen et Gauthier améliorant
$ \sqrt{3}/4 $, la supérieure une expression explicite en valeurs de la fonction Gamma trouvée par
Ahlfors et Grunsky en 1937 --- qui ont conjecturé que leur borne est la valeur exacte, conjecture
toujours indécise. La constante de Landau, sa compagne, est dans le même état. André Bloch a
démontré le théorème en 1925 alors qu'il était interné dans un hôpital psychiatrique, où il a fait
toutes ses mathématiques et correspondu avec Hadamard et Pólya ; la constante qui porte son nom
résiste à l'évaluation exacte depuis un siècle.

\begin{refs}
\item Contexte : Théorème de Bloch (analyse complexe) --- Wikipédia (en anglais) (\url{https://en.wikipedia.org/wiki/Bloch%27s_theorem_(complex_analysis)})
\item Lecture : Conway, \textit{Functions of One Complex Variable I}, ch. XII
\end{refs}

\bigskip

\begin{problem}[{La conjecture de la valeur moyenne de Smale --- Recherche, short}]
\textbf{CA-30.} \textbf{Problème.} Soit $ p $ un polynôme complexe de degré $ d \ge 2 $
avec $ p'(0) \ne 0 $. Démontrez ou réfutez qu'un certain point critique $ \zeta $ de $ p $
vérifie
\[ \left| \frac{p(\zeta) - p(0)}{\zeta \, p'(0)} \right| \;\le\; \frac{d - 1}{d}. \]
\end{problem}

\noindent \textbf{Statut : ouvert depuis 1981.} Smale a soulevé la question en analysant la
méthode de Newton dans ses travaux sur la complexité de la résolution des équations polynomiales, et
a démontré l'inégalité avec la constante $ 4 $ à la place de $ (d-1)/d $ ; les améliorations
depuis sont marginales, et la constante conjecturée ne peut pas être abaissée, puisque
$ p(z) = z^d - \lambda z $ l'atteint exactement. Le problème se tient aux côtés de celui de Sendov
(CA-31) dans la petite famille des questions sur les cachettes possibles des points critiques d'un
polynôme --- élémentaires à énoncer, et obstinées hors de toute proportion.

\begin{refs}
\item Contexte : Problème de la valeur moyenne --- Wikipédia (en anglais) (\url{https://en.wikipedia.org/wiki/Mean_value_problem})
\item Contexte : Théorème de Gauss--Lucas --- Wikipédia (\url{https://fr.wikipedia.org/wiki/Th%C3%A9or%C3%A8me_de_Gauss-Lucas})
\item Article : S. Smale, \og{} The fundamental theorem of algebra and complexity theory \fg{} (1981), Bull. Amer. Math. Soc. 4
\end{refs}

\bigskip

\begin{problem}[{La conjecture de Sendov --- Recherche}]
\textbf{CA-31.} \textbf{Problème.} Soit $ p(z) = (z - r_1)(z - r_2)\cdots(z - r_n) $ avec $ n \ge 2 $ et toutes
les racines dans le disque unité fermé. La conjecture de Sendov affirme qu'à distance au plus
$ 1 $ de chaque racine $ r_k $ se trouve un point critique de $ p $ --- un zéro de
$ p' $.
\begin{enumerate}
\item[(a)] Démontrez le théorème de Gauss--Lucas : tout zéro de $ p' $ est dans l'enveloppe convexe des
zéros de $ p $. Déduisez-en que sous l'hypothèse ci-dessus tous les points critiques sont dans le
disque unité fermé, et donc que la conjecture est vraie sans réserve pour toute racine située à
l'origine.
\item[(b)] Démontrez la conjecture pour $ n = 2 $ par calcul direct, et examinez $ p(z) = z^n - 1 $,
dont tous les points critiques sont à l'origine, à distance exactement $ 1 $ de chaque racine --- de
sorte qu'aucun rayon plus petit que $ 1 $ ne conviendrait, et que la conjecture, si elle est vraie,
est optimale.
\item[(c)] Démontrez la conjecture pour $ n = 3 $ --- un calcul, mais un calcul réellement pénible --- puis
lisez le traitement des grands degrés par Tao en 2020 et son exposé de 2026 de la démonstration
complète, et rendez compte de la difficulté : pourquoi les configurations quasi extrémales rendent
l'uniformité en $ n $ si difficile à obtenir, et comment une démonstration n'utilisant guère plus
que le théorème fondamental de l'algèbre a pu rester cachée soixante ans.
\end{enumerate}
\textit{Statut : très récemment résolue, et le dossier n'est pas encore refermé. Blagovest Sendov a
proposé la conjecture vers 1958 (elle a circulé des années durant sous le nom de conjecture
d'Ilieff) ; elle a été vérifiée pour les degrés inférieurs à 9 par Brown et Xiang, et Tao l'a
démontrée en 2020 pour tous les degrés suffisamment grands, par un argument de compacité et de
prolongement analytique qui ne donnait aucun seuil explicite. En août 2026, Lech Mazur a annoncé une
démonstration assistée par ordinateur du cas général, obtenue avec un système d'IA et vérifiée dans
l'assistant de preuve Lean ; Tao l'a ensuite digérée, reformalisée et simplifiée, et a utilisé le
même argument pour régler la conjecture plus forte de Phelps--Rodriguez. À l'heure où ces lignes sont
écrites, le résultat n'est pas passé par l'expertise conventionnelle d'une revue --- c'est la
formalisation vérifiée par machine qui en tient lieu de garantie.}
\end{problem}

\noindent La conjecture est une question sur les cachettes possibles des points critiques
d'un polynôme, et c'est l'illustration la plus tranchante de l'écart entre ce qui est facile à
énoncer et ce qui est facile à démontrer : les exemples extrémaux sont des quasi-égalités partout,
si bien qu'aucun argument souple ne survit. Sa résolution en 2026 est aussi un jalon d'un autre
ordre. La première démonstration a été produite par un système d'IA et comptait 90 000 lignes de
Lean ; un mathématicien humain l'a ensuite lue, réorganisée en un argument n'utilisant que le
théorème fondamental de l'algèbre et l'inégalité de Maclaurin, et l'a étendue. Quoi qu'on pense de
la méthode, l'épistémologie de cette page --- les références, les démonstrations et le crayon du
lecteur --- doit désormais peser un nouveau genre de source.

\begin{refs}
\item Contexte : Conjecture d'Iliev--Sendov --- Wikipédia (\url{https://fr.wikipedia.org/wiki/Conjecture_d'Iliev-Sendov})
\item Contexte : Théorème de Gauss--Lucas --- Wikipédia (\url{https://fr.wikipedia.org/wiki/Th%C3%A9or%C3%A8me_de_Gauss-Lucas})
\item Article : T. Tao, \og{} Sendov's conjecture for sufficiently high degree polynomials \fg{} (2020), arXiv:2012.04125 (\url{https://arxiv.org/abs/2012.04125})
\item Lecture : T. Tao, \og{} A digestion of the proof of Sendov's conjecture \fg{} (billet de blog, 12 août 2026), terrytao.wordpress.com (\url{https://terrytao.wordpress.com/2026/08/12/a-digestion-of-the-proof-of-sendovs-conjecture/})
\end{refs}

\bigskip

\begin{problem}[{L'ensemble de Mandelbrot est-il localement connexe ? --- Recherche, short}]
\textbf{CA-40.} \textbf{Problème.} Soit $ M $ l'ensemble de Mandelbrot, l'ensemble des
paramètres $ c \in \mathbb{C} $ pour lesquels l'orbite de $ 0 $ sous $ z \mapsto z^2 + c $
reste bornée ; il est compact, et Douady et Hubbard ont démontré qu'il est connexe. Déterminez si
$ M $ est \textit{localement connexe} --- si tout point de $ M $ possède des voisinages relatifs
connexes arbitrairement petits.
\end{problem}

\noindent \textbf{Statut : ouvert.} Connue sous le nom de MLC, la conjecture est le problème
ouvert central de la dynamique complexe en dimension un. Douady et Hubbard ont démontré la connexité
au début des années 1980 en construisant un isomorphisme conforme explicite entre le complémentaire
de $ M $ et le complémentaire du disque unité fermé --- de l'application conforme pure, dans
l'esprit de CA-11 et CA-13 --- et ils ont montré que la connexité locale fournirait un modèle
combinatoire complet de $ M $ en \og{} disque pincé \fg{}, et avec lui la conjecture d'hyperbolicité,
selon laquelle les paramètres hyperboliques sont denses dans la famille quadratique. Yoccoz a
démontré MLC en tout paramètre finiment renormalisable, travail récompensé par une médaille Fields
en 1994, et Lyubich et d'autres ont progressé dans le domaine infiniment renormalisable ; le cas
général reste ouvert à l'heure où ces lignes sont écrites.

\begin{refs}
\item Contexte : Ensemble de Mandelbrot --- Wikipédia (\url{https://fr.wikipedia.org/wiki/Ensemble_de_Mandelbrot})
\item Contexte : Conjecture de Fatou (densité de l'hyperbolicité) --- Wikipédia (en anglais) (\url{https://en.wikipedia.org/wiki/Fatou_conjecture}), Ensemble de Julia --- Wikipédia (\url{https://fr.wikipedia.org/wiki/Ensemble_de_Julia})
\item Lecture : Milnor, \textit{Dynamics in One Complex Variable}
\item Lecture : Carleson \& Gamelin, \textit{Complex Dynamics} (Springer, 1993)
\end{refs}

\bigskip

\begin{problem}[{Le problème de la couronne --- Recherche}]
\textbf{CA-41.} \textbf{Problème.} Soit $ H^{\infty} $ l'algèbre des fonctions holomorphes bornées sur le disque
unité $ \mathbb{D} $, munie de la norme sup. Le théorème de la couronne de Carleson (1962) dit :
si $ f_1, \dots, f_n \in H^{\infty} $ et s'il existe un $ \delta > 0 $ tel que
$ |f_1(z)| + \cdots + |f_n(z)| \ge \delta $ pour tout $ z \in \mathbb{D} $, alors il existe
$ g_1, \dots, g_n \in H^{\infty} $ tels que
\[ f_1 g_1 + \cdots + f_n g_n = 1 . \]
\begin{enumerate}
\item[(a)] Démontrez la moitié facile : si l'équation de Bézout a une solution dans $ H^{\infty} $, alors
la minoration $ \sum_j |f_j| \ge \delta $ doit valoir pour un certain $ \delta > 0 $. Réglez
complètement le cas $ n = 1 $.
\item[(b)] Démontrez que le théorème équivaut à l'énoncé selon lequel $ \mathbb{D} $ est dense dans
l'espace des idéaux maximaux de $ H^{\infty} $ --- l'énoncé qu'il n'y a pas de \og{} couronne \fg{} d'idéaux
supplémentaires autour du disque. (Identifiez les points de $ \mathbb{D} $ aux caractères
d'évaluation, et utilisez la correspondance entre idéaux maximaux et caractères d'une algèbre de
Banach commutative.)
\item[(c)] Construisez des exemples instructifs : montrez qu'une suite d'interpolation et son produit de
Blaschke (CA-35) fournissent des fonctions de petite minoration conjointe, et vérifiez à la main ce
qu'affirme le théorème pour $ f_1(z) = z $ et $ f_2 = $ un produit de Blaschke ne s'annulant
qu'en $ 0 $.
\item[(d)] Lisez la démonstration : l'argument original de Carleson via le théorème d'interpolation et les
mesures de Carleson, ou l'argument $ \bar\partial $ ultérieur de Wolff, et rendez compte de
l'endroit où la géométrie du disque intervient. Affrontez ensuite le problème ouvert : existe-t-il
un théorème de la couronne pour le polydisque ou pour la boule de $ \mathbb{C}^n $ avec
$ n \ge 2 $, ou pour un domaine plan quelconque ?
\end{enumerate}

\end{problem}

\noindent \textbf{Statut : résolu pour le disque ; ouvert en général.} Kakutani a demandé en
1941 si le disque pouvait avoir une couronne, et la question a résisté vingt ans d'efforts jusqu'à
ce que Lennart Carleson la règle en 1962 avec la machinerie --- suites d'interpolation, mesures de
Carleson --- qui a réorganisé toute la théorie des espaces de Hardy ; Thomas Wolff a trouvé une
démonstration bien plus courte en 1979 par des techniques $ \bar\partial $. Ce qui se passe au-delà
du disque est réellement inconnu : Cole a construit une surface de Riemann possédant une couronne,
donc le théorème n'est pas formel, et les cas du polydisque, de la boule en plusieurs variables et
d'un domaine plan général restent ouverts. C'est un problème ouvert inhabituel en ceci que son cas
résolu est l'une des grandes prouesses techniques de l'analyse du vingtième siècle, et que personne
ne sait si l'obstruction en dimension supérieure est une idée manquante ou un véritable
contre-exemple.

\begin{refs}
\item Contexte : Théorème de la couronne --- Wikipédia (en anglais) (\url{https://en.wikipedia.org/wiki/Corona_theorem})
\item Contexte : Espace de Hardy --- Wikipédia (\url{https://fr.wikipedia.org/wiki/Espace_de_Hardy}), Mesure de Carleson --- Wikipédia (\url{https://fr.wikipedia.org/wiki/Mesure_de_Carleson})
\item Lecture : Garnett, \textit{Bounded Analytic Functions}, ch. VIII
\item Article : L. Carleson, \og{} Interpolations by bounded analytic functions and the corona problem \fg{} (1962), \textit{Annals of Mathematics} 76
\end{refs}

\bigskip


\section*{Pour aller plus loin}


\vfill
\noindent\emph{Hors de la Caverne} --- des probl\`emes, pas de r\'eponses.
\end{document}
